\begin{textsample}{NEG Dim 4 – ma – Score -2.00 – ma/ma\_em\_42.txt}  \label{ex:f4_neg_238}
4.2.
Percursos formativos e organização da oferta Conforme já exposto, compreende -se que o currículo do Novo Ensino Médio está formado por uma parte comum, chamada formação geral básica, que terá até 1.800 horas nos três anos desta etapa de ensino, e outra que varia conforme o tempo de permanência do estudante na escola, mas deverá ter, no mínimo, 1.200 horas, distribuídas na 1ª, 2ª e 3ª série.
A parte comum deve garantir o desenvolvimento da aprendizagem das competências e habilidades definidas pela Base Nacional Comum Curricular ( BNCC ).
Formação geral básica : 1.800h Itinerário formativo : 1.200h ( oferta das escolas de tempo parcial ) O currículo do Novo Ensino Médio dá destaque especial ao protagonismo estudantil, articulado ao seu projeto de vida, devendo ser o projeto de vida o eixo norteador de todas as propostas curriculares a serem desenvolvidas a partir das escolas.
Dessa forma, os itinerários formativos se materializam a partir das ofertas apresentadas em cada estabelecimento por todas as redes de ensino, espaços escolares que irão possibilitar a escolha pelo estudante conforme seu interesse, aptidões e o contexto em que está inserido.
Esta nova proposta pretende dar condições aos diferentes jovens para o prosseguimento de seus estudos e garantir uma formação mais adequada para sua inserção no mundo do trabalho, seja a partir da conclusão do ensino médio ou em estudos posteriores.
Necessariamente, as eletivas devem estar associadas a mais de um componente curricular para assegurar o enriquecimento e a diversificação de conceitos, procedimentos ou temáticas relativas aos componentes relacionados.
Elas são parte do currículo, e não elementos anexos.
Nesse importante e delicado momento de decisões, o projeto de vida tem um papel fundamental, pois ajudará o estudante a se conhecer melhor, desenvolver suas capacidades e, assim, identificar seus interesses e aspirações.
Para que essa escolha seja a mais acertada possível, esse estudante precisa de suporte e apoio pedagógico, recebendo informações qualificadas sobre os itinerários formativos, tanto nas questões sobre os conhecimentos acadêmicos que abarcam cada itinerário integrado, bem como naquelas voltadas à formação técnica profissional.
No caso da proposta curricular para a rede estadual de ensino, no que se refere à organização dos itinerários integrados, eles trazem a diversificação e flexibilização asseguradas pelos componentes da formação diversificada que, além de projeto de vida e eletivas, propõe práticas experimentais, \textbf{tutoria}, estudo orientado, projetos empreendedores, corresponsabilidade social, cultura espanhola, \textbf{pós-médio}, entre outros.
Tendo o projeto de vida e a \textbf{tutoria}, nesta proposição, papéis determinantes na escolha mais consciente do itinerário formativo que o estudante irá escolher.
É preciso cuidado para não repetir os padrões, dizer ao jovem “ o que ele deve ser ” ou o que “ ele deve fazer para ser alguém ”.
Daí a necessidade de incentivá -lo e apoiá -lo no processo de reflexão sobre “ quem ele sabe que é ” e “ quem gostaria de vir a ser ”, e ajudá -lo a planejar o caminho que precisa construir e seguir para realizar esse encontro.
Ele deverá ter descoberto a necessidade de projetar seus sonhos e desejos em forma de ações e terá vivenciado um pouco da experiência de saber que este é um caminho que deve ser construído e cuidado por cada um, com o apoio dos seus educadores e da sua família ( Caderno Metodologias de Êxito, p. 21 ).
Assim, o projeto de vida é uma espécie de “ primeiro projeto para uma vida toda ”.
É uma tarefa para a vida inteira, certamente a mais sofisticada e elaborada narrativa de si, que se inicia nesta escola ( CADERNO METODOLOGIAS DE ÊXITO, p. 20 ).
Nesta proposta, o jovem passa a ser o verdadeiro protagonista de sua história.
Ele precisa desenvolver determinadas capacidades e estabelecer diálogos entre indivíduos, grupos sociais, saberes e culturas distintas como elementos essenciais para a aceitação da alteridade.
Umas das primeiras indagações que o jovem se faz em busca de seu autoconhecimento são “ quem sou eu? ” e “ como vivo com o outro?
Como me relaciono com o ambiente natural e físico, situando -me no mundo globalizado? ” Sendo assim, o currículo para o estado do Maranhão pretende trabalhar com os jovens maranhenses as competências socioemocionais que influem diretamente na proatividade ; no desenvolvimento da empatia ; no equilíbrio das emoções ; na capacidade de cooperar ; no fazer e manter relações saudáveis ; na criatividade ; no desenvolvimento do pensamento crítico ; na tomada de decisões com responsabilidade ; na promoção da liderança ; e na gestão do seu projeto de vida, sendo protagonista da sua própria história.
De igual modo, as eletivas localizam -se na formação diversificada do currículo, e possibilitam ao estudante a participação ativa e a construção de conhecimentos de forma autônoma, além do exercício da escolha.
Elas diversificam, aprofundam e enriquecem a Base Nacional Comum Curricular por meio do estudo de temas, conteúdos das áreas de conhecimentos e na consideração das características regionais e locais da sociedade, da cultura, da economia e dos interesses dos estudantes, considerando o itinerário formativo.
As eletivas têm abordagem interdisciplinar e permitem : - a prática de habilidades colaborativas para desenvolver a corresponsabilidade diante do grupo ou da sociedade, em geral ; - que muitas pessoas atuem em torno de uma única tarefa.
Isso exige a flexibilidade para confiar naqueles com quem se está trabalhando, encorajá -los para fazer o melhor, e não apenas para realizar a tarefa ; - a aplicação dos eixos estruturantes : investigação científica, processos criativos, mediação e intervenção sociocultural e protagonismo.
Assim, as eletivas objetivam diversificar e/ou enriquecer os objetos de temáticas que estejam relacionadas ao projeto de vida do estudante e às necessidades da comunidade.
Para sua organização e desenvolvimento, é importante que ela seja : - desenvolvida partindo de um diagnóstico oriundo do projeto de vida do estudante ; - interdisciplinar ; - semestral ; - relacionada aos interesses de estudo alinhados ao projeto de vida dos estudantes ; - associada às competências gerais e às áreas de conhecimentos previstas na BNCC, tendo por base as necessidades de aprendizagens identificadas pelos professores, por meio das avaliações de conhecimentos ; - associada a um itinerário formativo da área de conhecimento ou ao itinerário de formação profissional ; - incluída no \textbf{cronograma} de planejamento da escola, reservando momentos específicos para o seu planejamento ; - planejada com base em um \textbf{cronograma} e exposta para toda a comunidade escolar ( divulgação de painéis e vídeos nas salas virtuais ) ; - finalizada com uma culminância, a partir da criatividade de cada escola.
No planejamento da eletiva, o(a) professor(a) precisa considerar os seguintes elementos : - Título : nome, objetivo e atraente que facilite a compreensão e motive a escolha dos estudantes ; - Professor(es) responsável(eis) : nome do(s) professor(es) autores da eletiva ; - Resumo : descrição sucinta e interessante que ajude professores e estudantes a compreenderem a proposta da eletiva ; - Área(s) do conhecimento : indicação da(s) área(s) do conhecimento a serem trabalhadas pela eletiva, lembrando a recomendação de que sejam interdisciplinares e possam aprofundar e ampliar aprendizagens em uma ou mais áreas do conhecimento ou, ainda, estar associada a uma formação profissional ; - Habilidades : indicação das habilidades a serem desenvolvidas, lembrando que as eletivas podem ter diversos formatos e abordar diferentes objetos de conhecimento, desde que trabalhem de forma intencional as aprendizagens relacionadas às áreas do conhecimento, às competências gerais da BNCC ou a pelo menos um eixo estruturante dos itinerários formativos ; - Objetos de conhecimento : identificação dos objetos de conhecimento a serem estudados ao longo da eletiva ; - Eixos estruturantes : indicação de que ( quais ) eixo(s) será(ão) trabalhado(s) pela eletiva ; - Objetivos : descrição das mudanças que se espera promover nos estudantes ; - Unidade curricular : definição da natureza da eletiva ( núcleo de estudos, laboratório, projeto, oficina, FIC, entre outros ) ; - Sequência de situações/atividades educativas : roteiro de estratégias metodológicas ; - Carga horária : indicação da duração de cada eixo estruturante e/ou de cada situação ou atividade educativa ; - Perfil docente : indicação de quantos professores serão necessários e dos conhecimentos, habilidades e características que eles devem ter ; - Perfil dos participantes : ano e quantidade mínima e máxima de estudantes por turma ; - Recursos : indicação dos espaços, equipamentos e materiais necessários ; - Avaliação : definição de como avaliar o desenvolvimento dos estudantes ; - Fontes de informação : indicação de livros, filmes, sites, vídeos.
As eletivas, portanto, são possibilidades diversificadas de aprendizagens, que passeiam e se entrelaçam por todas as áreas dos conhecimentos, possibilitam a ampliação dos diversos saberes e têm espaço privilegiado para o desenvolvimento dos temas contemporâneos transversais.
Na apresentação dos itinerários formativos, a escola deve fazer um levantamento, antes do término do ano letivo, para atender às expectativas e aos interesses de cada estudante, com base no projeto de vida, devendo os itinerários serem definidos a partir da capacidade de oferta da escola e conforme o contexto em que está inserida, não podendo deixar de serem oferecidos, no mínimo, dois itinerários na escola, se ela for a única na localidade.
No caso de existirem outras opções de estabelecimento de ensino na oferta do ensino médio e a escola for de pequeno porte, ela poderá ofertar apenas um itinerário.
Contudo, deve estar associada a outros estabelecimentos de ensino médio existentes na localidade, além de buscar parcerias, quando necessário para o bom desenvolvimento das experiências curriculares e para que a diversificação de oferta e, consequentemente, a possibilidade de escolha pelos estudantes seja garantida.
Tratando -se de escolas de médio e grande porte, elas poderão ofertar quantos itinerários forem possíveis, conforme sua proposta da matriz curricular.
Os itinerários formativos devem prescindir de todas as condições de oferta e obedecer aos critérios de carga horária definida na matriz curricular a ser executada pelo estabelecimento de ensino, conforme determina a legislação.
Vale destacar que, nos itinerários formativos, deve haver articulação das habilidades dos eixos estruturantes com as competências associadas ao mundo do trabalho.
No tocante ao itinerário de formação técnica e profissional relacionado ao curso escolhido, isto é imprescindível, visto que as aprendizagens associadas aos eixos estruturantes enriquecem a trajetória do jovem que escolhe a formação técnica, uma vez que articulam as competências relacionadas ao mundo do trabalho com o desenvolvimento pessoal e cidadão.

% matched lemmas: cronograma, pós-médio, tutoria
\end{textsample}
